\section{Related Work}
\label{sec:related}

\textbf{NMS-free detectors.}
Set-prediction detectors remove non-maximum suppression (NMS) by
one-to-one matching~\cite{carion2020detr}; RT-DETR made this regime
real-time on GPUs~\cite{zhao2024rtdetr}. YOLOv10 introduced consistent dual
assignments so that a convolutional YOLO can be deployed with a
one-to-one head and no NMS~\cite{wang2024yolov10}, and YOLO26 continues
this design with an end-to-end head~\cite{jocher2026yolo26,ultralytics2026yolo26}. With NMS
gone, the remaining host work is decode, letterbox, and a short
top-$k$ row conversion. This makes the rest of the pipeline more visible,
and it is why we study these models.

\textbf{Compression for edge detection.}
Structured channel pruning~\cite{li2017pruning,fang2023depgraph},
quantization-aware training~\cite{jacob2018quantization} and knowledge
distillation~\cite{hinton2015distilling} are the standard tools for
making detectors cheaper. They are usually judged by parameter count,
FLOPs or isolated engine latency. We run these tools on the same
deployment stack and report them in end-to-end units. Our
negative results do not show that compression is ineffective in
general. They show that, for a nano-scale NMS-free model on this
device, the model is not where the time goes.

\textbf{Latency prediction, DVFS and thermal control.}
nn-Meter showed that kernel-level effects make FLOPs a poor latency
proxy on edge devices~\cite{zhang2021nnmeter}. Tang et al.\ measured how
GPU core and memory frequencies change the time and energy of deep
learning workloads, and found that the effect differs across
networks~\cite{tang2019gpudvfs}. Nabavinejad et al.\ chose batch size and
GPU frequency jointly for inference under latency limits~\cite{nabavinejad2022batching}.
Other DVFS work on mobile and edge processors learns frequency policies
that avoid throttling~\cite{kim2021ztt}, models GPU latency under
DVFS~\cite{han2025dvfs}, or controls frequency per operator on
Jetson-class devices~\cite{zhang2026sparsedvfs}. SparseDVFS also
documents, as motivation, the antagonistic effect of the stock CPU and
GPU governors on a Jetson Orin Nano: the CPU governor down-clocks the
host during GPU-heavy phases, so the host is slow exactly when the GPU
next requests data~\cite{zhang2026sparsedvfs}. \Cref{sec:sched}
measures the same coupling inside a detection pipeline and traces it
to the CUDA wait policy. Work on 3D mobile games
showed that managing CPU and GPU frequencies jointly saves power compared
with independent governors, because the two processors depend on each
other's progress~\cite{pathania2014cpugpu}. Closest to our setting,
Karzhaubayeva et al.\ characterize CNN workloads on a Jetson TX2 and
design integrated CPU--GPU DVFS governors for them~\cite{karzhaubayeva2023cpugpu},
Lotus manages CPU and GPU frequencies jointly for object detectors on a
Jetson Orin Nano to stabilize latency under thermal
variation~\cite{gong2024lotus}, and Kang shows that latency estimators
for deadline-aware governors on Jetson Orin must condition on the memory
clock~\cite{kang2026emc}. All of these works set
frequencies, or batch sizes, from outside the application --- they design
or evaluate frequency-control policies. We leave the
stock governors in place and ask how an application's duty cycle
interacts with them: the schedule changes the operating state the stock
governors observe, and thereby the measured cost of an otherwise fixed
engine. We find a similar CPU--GPU governor coupling in a
detection pipeline, and we test one standard kernel mechanism, utilization
clamping~\cite{linux2024uclamp}, as a remedy that needed no root privileges on our system.

\textbf{Characterization of edge inference.}
Mittal surveyed optimized deep learning on Jetson devices, covering
precision, framework and power-mode choices~\cite{mittal2019jetson}.
Hadidi et al.\ compared frameworks and devices for edge inference and
reported that, for the same model, the choice of framework changes
inference time and energy considerably~\cite{hadidi2019edge}. These studies
measure the model, or the framework that runs it. We measure the
application around it: we hold the engine fixed and change only host
code and governor settings.

\textbf{End-to-end costs outside the model.}
The ``AI tax'' study showed that pre- and post-processing can dominate
accelerated AI applications~\cite{richins2020aitax}. On Jetson, JPEG
decode and preprocessing have been offloaded to dedicated
units~\cite{chen2026nvjpeg}, training pipelines have been characterized
for I/O pipelining and power modes~\cite{prashanthi2023jetson}, and the
latency distribution of a YOLO pipeline on Jetson Orin NX has been
modeled stage by stage as a function of image resolution~\cite{zhang2026resolution}.
MLPerf Inference standardizes measurement scenarios, separating offline
throughput from server and single-stream latency~\cite{reddi2020mlperf}.
For autonomous-driving detection, R-TOD analyzes and minimizes the
end-to-end delay from frame exposure to reported detections by
restructuring capture and pipelining (on-demand capture, zero-slack and
contention-free pipelines) without changing the
network~\cite{jang2020rtod}; our arrival-aware variant
(\cref{sec:stream}) is a simple member of the same family, and we
measure its interaction with the stock governors and the idle floor
rather than worst-case delay. We
do not claim that pipelining or preprocessing bottlenecks are new. What we
add is a measured account, on one device and stack, of how the
governor's response to the application duty cycle changes the
apparent cost of the \emph{model}. We also show how this reorders
deployment decisions such as compression and resolution.

\textbf{Adaptive resolution.}
Choosing input scale per image or frame can save compute in
detection~\cite{chin2019adascale,zhu2021drnet}. Our frozen pre-inference
router is a simple member of this family. We report it because it
failed held-out verification, and because its expected saving is smaller
than the pipeline effects we measure.

\begin{table}[t]
\centering
\caption{Positioning against the closest measurement studies. ``Freq.''
is how frequency enters the experiment: set directly by the
experimenter, replaced by a newly designed governor (``new gov.''),
or left to the stock governor and observed (``modif.'': a modified
stock governor; ``ucl.'': per-task utilization clamping). ``Pipe.'':
the whole decode-to-detections pipeline is measured, not only the
model. ``Acc.'': accuracy is held constant or verified per variant.
``Stream.'': camera-rate arrivals are measured, not only backlog
throughput. SLO: service-level objective. A dash means the work does
not address the aspect.}
\label{tab:related}
\scriptsize
\setlength{\tabcolsep}{1.6pt}
\begin{tabular}{@{}l l cccc@{}}
\toprule
Work & Workload & Freq. & Pipe. & Acc. & Stream. \\
\midrule
Tang et al.~\cite{tang2019gpudvfs} & DNN time/energy & direct & -- & -- & -- \\
Nabavinejad et al.~\cite{nabavinejad2022batching} & batching w/ SLOs & direct & -- & -- & -- \\
SparseDVFS~\cite{zhang2026sparsedvfs} & per-operator DVFS & direct & -- & yes & -- \\
Pathania et al.~\cite{pathania2014cpugpu} & 3D mobile games & modif. & partial & -- & -- \\
Karzhaubayeva~\cite{karzhaubayeva2023cpugpu} & CNNs on TX2 & new gov. & -- & -- & -- \\
Lotus~\cite{gong2024lotus} & detec., Orin Nano & new gov. & -- & -- & -- \\
Kang~\cite{kang2026emc} & gov.\ latency est. & direct & -- & -- & -- \\
AI tax~\cite{richins2020aitax} & pre/post-processing & -- & yes & -- & -- \\
Stage-wise latency~\cite{zhang2026resolution} & YOLO on Orin NX & stock & yes & -- & -- \\
This work & detection, 3 boards & stock${+}$ucl. & yes & yes & yes \\
\bottomrule
\end{tabular}
\end{table}

\Cref{tab:related} summarizes the positioning. The gap we fill is not
pipelining, preprocessing or DVFS per se, but the interaction with the
\emph{stock} governors under the application's own schedule and arrival
process, with accuracy held exactly constant so that the timing effects
are attributable.
